% D1-S12 Chimera
% Auto-generated from markdown source
% origin: ai

\documentclass[12pt, a4paper]{article}
\usepackage[utf8]{inputenc}
\usepackage[T1]{fontenc}
\usepackage{amsmath, amssymb, amsfonts, amsthm}
\usepackage{textcomp}
\usepackage{geometry}
\usepackage{hyperref}
\usepackage{booktabs}
\usepackage{graphicx}
\usepackage{fancyhdr}
\usepackage{tcolorbox}
\usepackage{enumitem}

\geometry{margin=1in}
\pagestyle{fancy}
\fancyhf{}
\rhead{D1-S12 Chimera}
\lhead{\thepage}

\newtheorem{definition}{Definition}[section]
\newtheorem{theorem}{Theorem}[section]
\newtheorem{corollary}{Corollary}[section]
\newtheorem{proposition}{Proposition}[section]
\newtheorem{lemma}{Lemma}[section]
\newtheorem{conjecture}{Conjecture}[section]

\newcommand{\calI}{\mathcal{I}}
\newcommand{\calP}{\mathcal{P}}
\newcommand{\calE}{\mathcal{E}}
\newcommand{\calS}{\mathcal{S}}
\newcommand{\calH}{\mathcal{H}}
\newcommand{\calO}{\mathcal{O}}
\newcommand{\calF}{\mathcal{F}}
\newcommand{\calU}{\mathcal{U}}
\newcommand{\calV}{\mathcal{V}}
\newcommand{\calA}{\mathcal{A}}
\newcommand{\calG}{\mathcal{G}}
\newcommand{\calB}{\mathcal{B}}
\newcommand{\calC}{\mathcal{C}}
\newcommand{\calL}{\mathcal{L}}
\newcommand{\calM}{\mathcal{M}}
\newcommand{\calX}{\mathcal{X}}
\newcommand{\calY}{\mathcal{Y}}
\newcommand{\calW}{\mathcal{W}}
\newcommand{\calT}{\mathcal{T}}
\newcommand{\calZ}{\mathcal{Z}}
\newcommand{\calR}{\mathcal{R}}
\newcommand{\calQ}{\mathcal{Q}}
\newcommand{\calD}{\mathcal{D}}
\newcommand{\calN}{\mathcal{N}}
\newcommand{\calK}{\mathcal{K}}

\newcommand{\Tr}{\text{Tr}}
\newcommand{\E}{\mathbb{E}}
\newcommand{\diff}{\mathrm{d}}
\newcommand{\KL}{D_{\mathrm{KL}}}
\newcommand{\argmax}{\arg\max}
\newcommand{\argmin}{\arg\min}
\newcommand{\hbarF}{\hbar}

\newtcolorbox{thmbox}[1]{
  colback=blue!5!white, colframe=blue!75!black, fonttitle=\bfseries,
  title={#1}, breakable
}
\newtcolorbox{defbox}[1]{
  colback=green!5!white, colframe=green!75!black, fonttitle=\bfseries,
  title={#1}, breakable
}

\begin{document}

\title{D1-S12 Chimera}
\author{RALPH Loop}
\date{}
\maketitle

\begin{abstract}
This paper presents D1-S12 Chimera, formalizing the underlying principles using information-theoretic methods.

\textbf{Keywords}: D1-S12 Chimera, information theory, formalization.
\end{abstract}

\tableofcontents
\section{D1-S12 Chimera}


****D1-S12-r01chimera-cognitive-dialectic-theorem
****D9-IVD1-r02D4-S18-r23chimera S12-r02


\#\#

\textbf{ D4-S18-r23 N=256  chimera  gradient\_locking chimera}

D1-r02
\begin{itemize}
\item chimera (M=2)  /
\item gradient\_locking
\end{itemize}





\subsection{ N}


\textbf{ D1-S12-T1}

N

\[N \geq N_c \approx 500\]

\textbf{ sketch}

1. \textbf{D4-S18-r19}  chimera  N  500 +  K(x)

2. \textbf{D4-S18-r23 }
\begin{itemize}
\item N=256 + 4/4  gradient\_lockingMPI=0.00
\item chimera+SCI=0.496
\end{itemize}


3. \textbf{D1-r02 }
\begin{itemize}
\item chimera (M=2)  coherent  = incoherent  =
\item gradient\_locking
\end{itemize}


4. **** N < N\_c


\#\#

\subsubsection{ 1}


N\_c
-
-
\begin{itemize}
\item "" gradient\_locking
\end{itemize}


\subsubsection{ 2}


N
\begin{itemize}
\item N    chimera-like
\item N    global\_sync gradient\_locking
\end{itemize}



\#\#


1. **** N < 500  chimera gradient\_locking
2. ****


\subsection{ D9-IV}


D9-IV  (f) = (D, L, T)   T

\textbf{D1-S12-T1  T } =  spatially-heterogeneous

N < N\_c


* D4-S18-r24 *

\begin{thebibliography}{99}
\bibitem{ref1} Shannon, C.E. (1948). A Mathematical Theory of Communication.
\bibitem{ref2} Cover, T.M. (2006). Elements of Information Theory.
\end{thebibliography}

\end{document}